\olchapter{fol}{com}{The Completeness Theorem}



\olfileid{fol}{com}{int}

\olsection{Introduction}

The completeness theorem is one of the most fundamental results about
logic.  It comes in two formulations, the equivalence of which we'll
prove.  In its first formulation it says something fundamental about
the relationship between semantic consequence and our !!{derivation}
system: if !!a{sentence}~$!A$ follows from some !!{sentence}s
$\Gamma$, then there is also !!a{derivation} that establishes $\Gamma
\Proves !A$. Thus, the !!{derivation} system is as strong as it can
possibly be without proving things that don't actually follow.

In its second formulation, it can be stated as a model existence
result: every consistent set of !!{sentence}s is satisfiable.
Consistency is a proof-theoretic notion: it says that our
!!{derivation} system is unable to produce certain !!{derivation}s.
But who's to say that just because there are no !!{derivation}s of a
certain sort from~$\Gamma$, it's guaranteed that there is
!!a{structure}~$\Struct{M}$? Before the
completeness theorem was first proved---in fact before we had the
!!{derivation} systems we now do---the great German mathematician
David Hilbert held the view that consistency of mathematical theories
guarantees the existence of the objects they are about. He put it as
follows in a letter to Gottlob Frege:
\begin{quote}
  If the arbitrarily given axioms do not contradict one another with
  all their consequences, then they are true and the things defined by
  the axioms exist. This is for me the criterion of truth and
  existence. 
\end{quote}
Frege vehemently disagreed. The second formulation of the completeness
theorem shows that Hilbert was right in at least the sense that if the
axioms are consistent, then \emph{some}
!!{structure} exists that makes them all
true.

These aren't the only reasons the completeness theorem---or rather,
its proof---is important.  It has a number of important consequences,
some of which we'll discuss separately.  For instance, since any
!!{derivation} that shows $\Gamma \Proves !A$ is finite and so can
only use finitely many of the !!{sentence}s in~$\Gamma$, it follows by
the completeness theorem that if $!A$ is a consequence of~$\Gamma$, it
is already a consequence of a finite subset of~$\Gamma$.  This is
called \emph{compactness}.  Equivalently, if every finite subset of
$\Gamma$ is consistent, then $\Gamma$ itself must be consistent.

Although the compactness theorem follows from the completeness theorem
via the detour through !!{derivation}s, it is also possible to use the
\emph{the proof of} the completeness theorem to establish it
directly. For what the proof does is take a set of !!{sentence}s with
a certain property---consistency---and constructs !!a{structure} out
of this set that has certain properties (in this case, that it
satisfies the set). Almost the very same construction can be used to
directly establish compactness, by starting from ``finitely
satisfiable'' sets of !!{sentence}s instead of consistent ones.
The construction also yields other consequences, e.g.,
  that any satisfiable set of !!{sentence}s has a finite or
  !!{denumerable}s model.  (This result is called the
  L\"owenheim--Skolem theorem.)  In general, the construction of
  !!{structure}s from sets of !!{sentence}s is used often in logic,
  and sometimes even in philosophy.





\olfileid{fol}{com}{out}

\olsection{Outline of the Proof}

The proof of the completeness theorem is a bit complex, and upon first
reading it, it is easy to get lost.  So let us outline the proof.  The
first step is a shift of perspective, that allows us to see a route to
a proof.  When completeness is thought of as ``whenever $\Gamma
\Entails !A$ then $\Gamma \Proves !A$,'' it may be hard to even come
up with an idea: for to show that $\Gamma \Proves !A$ we have to find
!!a{derivation}, and it does not look like the hypothesis that
$\Gamma \Entails !A$ helps us for this in any way.  For some proof
systems it is possible to directly construct !!a{derivation}, but we
will take a slightly different approach.  The shift in perspective
required is this: completeness can also be formulated as: ``if
$\Gamma$ is consistent, it is satisfiable.''  Perhaps we can use the
information in~$\Gamma$ together with the hypothesis that it is
consistent to construct !!a{structure}
that satisfies every !!{sentence}
in~$\Gamma$.  After all, we know
what kind of !!{structure} we are looking
for: one that is as $\Gamma$ describes it!{}

If $\Gamma$ contains only atomic
  !!{sentence}s, it is easy to construct
a model for it. Suppose the atomic !!{sentence}s are all
  of the form $\Atom{P}{a_1,\dots,a_n}$ where the $a_i$ are
  !!{constant}s. All we have to do is come up with

  !!a{domain}~$\Domain{M}$ and an assignment for~$P$ so that
  $\Sat{M}{\Atom{P}{a_1,\ldots,a_n}}$. But that's not very hard:
  put $\Domain{M} = \Nat$, $\Assign{\Obj c_i}{M} = i$, and for every
  $\Atom{P}{a_1,\ldots,a_n} \in \Gamma$, put the tuple $\tuple{k_1,
    \dots, k_n}$ into $\Assign{P}{M}$, where $k_i$ is the index of the
  constant symbol~$a_i$ (i.e., $a_i \ident \Obj c_{k_i}$).

Now suppose $\Gamma$ contains some !!{formula}~$\lnot !B$, with $!B$
atomic.  We might worry that the construction of
$\Struct{M}$ interferes with the possibility
of making $\lnot !B$ true.  But here's where the consistency
of~$\Gamma$ comes in: if $\lnot !B \in \Gamma$, then $!B \notin
\Gamma$, or else $\Gamma$ would be inconsistent.  And if $!B \notin
\Gamma$, then according to our construction
of~$\Struct{M}$,
$\Sat/{M}{!B}$, so
$\Sat{M}{\lnot !B}$.  So far so good.

What if $\Gamma$ contains complex, non-atomic formulas? Say it
contains $!A \land !B$. To make that true, we should proceed as if
both $!A$ and $!B$ were in~$\Gamma$.  And if $!A \lor !B \in \Gamma$,
then we will have to make at least one of them true, i.e., proceed as
if one of them was in~$\Gamma$.

This suggests the following idea: we add additional !!{formula}s
to~$\Gamma$ so as to (a)~keep the resulting set consistent and
(b)~make sure that for every possible atomic !!{sentence}~$!A$, either
$!A$ is in the resulting set, or $\lnot !A$ is, and (c)~such that,
whenever $!A \land !B$ is in the set, so are both $!A$ and $!B$, if
$!A \lor !B$ is in the set, at least one of $!A$ or $!B$ is also, etc.
We keep doing this (potentially forever).  Call the set of all
!!{formula}s so added~$\Gamma^*$.  Then our construction above would
provide us with !!a{structure}~$\Struct{M}$
for which we could prove, by induction, that it satisfies all sentences
in~$\Gamma^*$, and hence also all sentence in~$\Gamma$
since $\Gamma \subseteq \Gamma^*$.  It turns out that guaranteeing (a)
and~(b) is enough. A set of sentences for which (b) holds is called
\emph{complete}. So our task will be to extend the consistent
set~$\Gamma$ to a consistent and complete set~$\Gamma^*$.


There is one wrinkle in this plan: if $\lexists[x][!A(x)] \in \Gamma$
we would hope to be able to pick some !!{constant}~$c$ and add $!A(c)$
in this process.  But how do we know we can always do that?  Perhaps we
only have a few !!{constant}s in our language, and for each one of
them we have $\lnot !A(c) \in \Gamma$.  We can't also add $!A(c)$,
since this would make the set inconsistent, and we wouldn't know
whether $\Struct{M}$ has to make $!A(c)$ or $\lnot !A(c)$ true.
Moreover, it might happen that $\Gamma$ contains only sentences in a
language that has no constant symbols at all (e.g., the language of
set theory).

The solution to this problem is to simply add infinitely many
constants at the beginning, plus sentences that connect them with the
quantifiers in the right way.  (Of course, we have to verify that this
cannot introduce an inconsistency.)

Our original construction works well if we only have !!{constant}s in
the atomic sentences.  But the language might also contain
!!{function}s.  In that case, it might be tricky to find the right
functions on~$\Nat$ to assign to these !!{function}s to make
everything work.  So here's another trick: instead of using $i$ to
interpret $\Obj c_i$, just take the set of !!{constant}s itself as the
domain.  Then $\Struct M$ can assign every !!{constant} to itself:
$\Assign{\Obj c_i}{M} = \Obj c_i$.  But why not go all the way: let
$\Domain{M}$ be all \emph{terms} of the language!{} If we do this,
there is an obvious assignment of functions (that take terms as
arguments and have terms as values) to !!{function}s: we assign to the
!!{function}~$\Obj f^n_i$ the function which, given $n$ terms $t_1$,
\dots,~$t_n$ as input, produces the term $\Obj f^n_i(t_1, \dots, t_n)$
as value.

The last piece of the puzzle is what to do with~$\eq$.  The
!!{predicate}~$\eq$ has a fixed interpretation: $\Sat{M}{\eq[t][t']}$
iff $\Value{t}{M} = \Value{t'}{M}$.  Now if we set things up so that the
!!{value} of a term~$t$ is $t$ itself, then this !!{structure}
will make \emph{no} sentence of the form $\eq[t][t']$ true
unless $t$ and $t'$ are one and the same term.  And of
course this is a problem, since basically every interesting theory in
a language with !!{function}s will have as theorems sentences
$\eq[t][t']$ where $t$ and $t'$ are not the same term (e.g., in
theories of arithmetic: $\eq[(\Obj 0+ \Obj 0)][\Obj 0]$).  To solve
this problem, we change the domain of~$\Struct M$: instead of using terms as
the objects in~$\Domain{M}$, we use sets of terms, and each set is so
that it contains all those terms which the sentences in~$\Gamma$
require to be equal.  So, e.g., if $\Gamma$ is a theory of arithmetic,
one of these sets will contain: $\Obj 0$, $(\Obj 0 + \Obj 0)$, $(\Obj
0 \times \Obj 0)$, etc.  This will be the set we assign to $\Obj 0$,
and it will turn out that this set is also the value of all the terms
in it, e.g., also of $(\Obj 0 + \Obj 0)$.  Therefore, the sentence
$\eq[(\Obj 0+ \Obj 0)][\Obj 0]$ will be true in this revised
!!{structure}.

So here's what we'll do. First we investigate the properties of
!!{complete} consistent sets, in particular we prove that
!!a{complete} consistent set contains $!A \land !B$ iff it contains
both $!A$ and~$!B$, $!A \lor !B$ iff it contains at least one of them,
etc. (\olref[ccs]{prop:ccs}). Then we define and
  investigate ``saturated'' sets of sentences. A saturated set is one
  which contains conditionals that link each quantified !!{sentence}
  to instances of it (\olref[hen]{defn:henkin-exp}). We show that any
  consistent set~$\Gamma$ can always be extended to a saturated
  set~$\Gamma'$ (\olref[hen]{lem:henkin}). If a set is consistent,
  saturated, and !!{complete} it also has the property that it
  contains $\lexists[x][!A(x)]$ iff it contains $!A(t)$
    for some closed term~$t$ and
  $\lforall[x][!A(x)]$ iff it contains $!A(t)$ for all
    closed terms~$t$ (\olref[hen]{prop:saturated-instances}).
We'll then take the saturated consistent
set~$\Gamma'$ and show that it can be extended
to a saturated, consistent, and
!!{complete} set~$\Gamma^*$ (\olref[lin]{lem:lindenbaum}). This set
$\Gamma^*$ is what we'll use to define our term
  model~$\Struct M(\Gamma^*)$.
The term model has the set of closed terms as its domain,
  and the interpretation of its !!{predicate}s is given by the atomic
  !!{sentence}s in~$\Gamma^*$ (\olref[mod]{defn:termmodel}).  We'll
use the properties of saturated, complete consistent
sets to show that indeed
$\Sat{M(\Gamma^*)}{!A}$ iff $!A \in
\Gamma^*$ (\olref[mod]{lem:truth}), and thus in particular,
$\Sat{M(\Gamma^*)}{\Gamma}$.
Finally, we'll consider how to define a term model if
  $\Gamma$ contains~$\eq$ as well
  (\olref[ide]{defn:term-model-factor}) and show that it
  satisfies~$\Gamma^*$ (\olref[ide]{lem:truth}).





\olfileid{fol}{com}{ccs}
      
\olsection{Complete Consistent Sets of \usetoken{P}{sentence}}

\begin{defn}[Complete set]
\ollabel{def:complete-set} A set~$\Gamma$ of !!{sentence}s is
\emph{!!{complete}} iff for any !!{sentence}~$!A$, either $!A \in
\Gamma$ or $\lnot !A \in \Gamma$.
\end{defn}

\begin{explain}
!!^{complete} sets of sentences leave no questions unanswered. For
any !!{sentence}~$!A$, $\Gamma$ ``says'' if $!A$ is true or false.  The
importance of !!{complete} sets extends beyond the proof of the
completeness theorem. A theory which is !!{complete} and
axiomatizable, for instance, is always decidable.
\end{explain}

\begin{explain}
!!^{complete} consistent sets are important in the completeness proof
since we can guarantee that every consistent set of
!!{sentence}s~$\Gamma$ is contained in !!a{complete} consistent
set~$\Gamma^*$.  !!^a{complete} consistent set contains, for each
!!{sentence}~$!A$, either $!A$ or its negation $\lnot !A$, but not
both. This is true in particular for atomic
!!{sentence}s, so from !!a{complete}
consistent set in a language suitably expanded by
!!{constant}s, we can construct
!!a{structure} where the
interpretation of !!{predicate}s is defined according to which
atomic !!{sentence}s are
in~$\Gamma^*$. This !!{structure} can then
be shown to make all !!{sentence}s in~$\Gamma^*$ (and hence also all
those in~$\Gamma$) true. The proof of this latter fact requires that
$\lnot !A \in \Gamma^*$ iff $!A \notin \Gamma^*$, $(!A \lor !B) \in
\Gamma^*$ iff $!A \in \Gamma^*$ or $!B \in \Gamma^*$, etc.
\end{explain}

In what follows, we will often tacitly use the properties of
reflexivity, monotonicity, and transitivity of $\Proves$ (see

  \tagrefs{prfSC/{fol:seq:ptn:sec},prfND/{fol:ntd:ptn:sec},prfAX/{fol:axd:ptn:sec},prfTab/{fol:tab:ptn:sec}}).

\begin{prop}
\ollabel{prop:ccs}
Suppose $\Gamma$ is !!{complete} and consistent. Then:
\begin{enumerate}
\item \ollabel{prop:ccs-prov-in} If $\Gamma \Proves !A$, then $!A \in
  \Gamma$.

\ollabel{prop:ccs-and} $!A \land !B \in \Gamma$
  iff both $!A \in \Gamma$ and $!B \in \Gamma$.

\ollabel{prop:ccs-or} $!A \lor !B \in \Gamma$ iff
  either $!A \in \Gamma$ or $!B \in \Gamma$.

\ollabel{prop:ccs-if} $!A \lif !B \in \Gamma$ iff
  either $!A \notin \Gamma$ or $!B \in \Gamma$.
\end{enumerate}
\end{prop}

\begin{proof}
Let us suppose for all of the following that $\Gamma$ is !!{complete} and
consistent.
\begin{enumerate}
\item If $\Gamma \Proves !A$, then $!A \in \Gamma$.

Suppose that $\Gamma \Proves !A$. Suppose to the contrary that $!A
\notin \Gamma$. Since $\Gamma$ is !!{complete}, $\lnot !A \in \Gamma$.
By

  \tagrefs{prfSC/{fol:seq:prv:prop:explicit-inc},
    prfND/{fol:ntd:prv:prop:explicit-inc},
    prfAX/{fol:axd:prv:prop:explicit-inc},
    prfTab/{fol:tab:prv:prop:explicit-inc}},
$\Gamma$ is inconsistent. This contradicts the assumption that
$\Gamma$ is consistent. Hence, it cannot be the case that $!A \notin
\Gamma$, so $!A \in \Gamma$.


Exercise.



First we show that if $!A \lor !B \in \Gamma$, then either $!A \in
\Gamma$ or $!B \in \Gamma$. Suppose $!A \lor !B \in \Gamma$ but $!A
\notin \Gamma$ and $!B \notin \Gamma$.  Since $\Gamma$ is
!!{complete}, $\lnot !A \in \Gamma$ and $\lnot !B \in \Gamma$. By

  \tagrefs{prfSC/{fol:seq:ppr:prop:provability-lor},
    prfND/{fol:ntd:ppr:prop:provability-lor},
    prfAX/{fol:axd:ppr:prop:provability-lor},
    prfTab/{fol:tab:ppr:prop:provability-lor}},
item (1), $\Gamma$ is inconsistent, a contradiction. Hence, either $!A
\in \Gamma$ or $!B \in \Gamma$.

For the reverse direction, suppose that $!A \in \Gamma$ or $!B \in
\Gamma$. By

  \tagrefs{prfSC/{fol:seq:ppr:prop:provability-lor},
    prfND/{fol:ntd:ppr:prop:provability-lor},
    prfAX/{fol:axd:ppr:prop:provability-lor},
    prfTab/{fol:tab:ppr:prop:provability-lor}}, item (2),
$\Gamma \Proves !A \lor !B$. By \olref{prop:ccs-prov-in}, $!A \lor
!B \in \Gamma$, as required.


Exercise.
\end{enumerate}
\end{proof}


\begin{prob}
Complete the proof of \olref[fol][com][ccs]{prop:ccs}.
\end{prob}







  

\olfileid{fol}{com}{hen}

\olsection{Henkin Expansion}

\begin{explain}
Part of the challenge in proving the completeness theorem is that the
model we construct from a complete consistent set~$\Gamma$ must make
all the quantified !!{formula}s in~$\Gamma$ true.  In order to
guarantee this, we use a trick due to Leon Henkin.  In essence, the
trick consists in expanding the language by infinitely many !!{constant}s
and adding, for each !!{formula} with one free !!{variable} $!A(x)$ a
formula of the form
$\lexists[x][!A(x)] \lif !A(c)$,
where $c$ is one of the new !!{constant}s.  When we construct the
!!{structure} satisfying~$\Gamma$, this will guarantee that each
true existential sentence has a witness
among the new constants.
\end{explain}

\begin{prop}
\ollabel{prop:lang-exp}
If $\Gamma$ is consistent in $\Lang L$ and $\Lang L'$ is obtained from
$\Lang L$ by adding !!a{denumerable} set of new !!{constant}s $\Obj d_0$,
$\Obj d_1$, \dots, then $\Gamma$ is consistent in~$\Lang L'$.
\end{prop}

\begin{defn}[Saturated set]
A set $\Gamma$ of !!{formula}s of a language $\Lang {L}$ is
\emph{saturated} iff for each !!{formula}~$!A(x) \in \Frm[L]$ with one
free !!{variable}~$x$ there is !!a{constant}~$c \in \Lang{L}$ such
that
$\lexists[x][!A(x)] \lif !A(c) \in \Gamma$.
\end{defn}

The following definition will be used in the proof of the next theorem.

\begin{defn}
\ollabel{defn:henkin-exp}
Let $\Lang L'$ be as in \olref{prop:lang-exp}.  Fix an enumeration
$!A_0(x_0)$, $!A_1(x_1)$, \dots of all !!{formula}s~$!A_i(x_i)$
of~$\Lang L'$ in which one variable ($x_i$) occurs free.  We define
the !!{sentence}s~$!D_n$ by induction on~$n$.

Let $c_0$ be the first !!{constant} among the $\Obj d_i$ we added
to~$\Lang{L}$ which does not occur in~$!A_0(x_0)$.  Assuming that
$!D_0$, \dots,~$!D_{n-1}$ have already been defined, let $c_n$ be the
first among the new !!{constant}s~$\Obj d_i$ that occurs neither in
$!D_0$, \dots,~$!D_{n-1}$ nor in~$!A_n(x_n)$.

Now let $!D_{n}$ be the !!{formula} $\lexists[x_{n}][!A_{n}(x_{n})] \lif
  !A_{n}(c_{n})$.
\end{defn}

\begin{lem}
\ollabel{lem:henkin}
Every consistent set~$\Gamma$ can be extended to a saturated
consistent set~$\Gamma'$.
\end{lem}

\begin{proof}
Given a consistent set of sentences~$\Gamma$ in a language~$\Lang{L}$,
expand the language by adding !!a{denumerable} set of new
!!{constant}s to form~$\Lang{L'}$.  By \olref{prop:lang-exp}, $\Gamma$
is still consistent in the richer language.  Further, let $!D_i$ be as
in \olref{defn:henkin-exp}.  Let
\begin{align*}
\Gamma_0 & = \Gamma \\
\Gamma_{n+1} & = \Gamma_n \cup \{!D_n \}
\end{align*}
i.e., $\Gamma_{n+1} = \Gamma \cup \{ !D_0, \dots, !D_n \}$, and let
$\Gamma' = \bigcup_{n} \Gamma_n$.  $\Gamma'$ is clearly saturated.

If $\Gamma'$ were inconsistent, then for some $n$, $\Gamma_n$ would be
inconsistent (Exercise: explain why).  So to show that $\Gamma'$ is
consistent it suffices to show, by induction on~$n$, that each
set~$\Gamma_n$ is consistent.

The induction basis is simply the claim that $\Gamma_0 = \Gamma$ is
consistent, which is the hypothesis of the theorem.  For the induction
step, suppose that $\Gamma_{n}$ is consistent but $\Gamma_{n+1} =
\Gamma_n \cup \{!D_n\}$ is inconsistent.  Recall that $!D_n$~is
$\lexists[x_{n}][!A_{n}(x_n)] \lif !A_{n}(c_{n})$,
where $!A_n(x_n)$ is !!a{formula} of $\Lang{L'}$ with only the
variable~$x_n$ free. By the way we've chosen the~$c_n$ (see
\olref{defn:henkin-exp}), $c_n$ does not occur in~$!A_n(x_n)$ nor
in~$\Gamma_n$.

If $\Gamma_n \cup \{!D_n\}$ is inconsistent, then $\Gamma_n
\Proves \lnot !D_n$, and hence both of the following hold:

\[
\Gamma_n \Proves \lexists[x_n][!A_n(x_n)]
\qquad
\Gamma_n \Proves \lnot !A_n(c_n)
\]
Since $c_n$ does not occur in
$\Gamma_n$ or in~$!A_n(x_n)$, 
\tagrefs{prfAX/{fol:axd:qpr:thm:strong-generalization},prfSC/{fol:seq:qpr:thm:strong-generalization},prfND/{fol:ntd:qpr:thm:strong-generalization},prfTab/{fol:tab:qpr:thm:strong-generalization}} applies.
From $\Gamma_n \Proves \lnot !A_n(c_n)$,
we obtain
$\Gamma_n \Proves \lforall[x_n][\lnot !A_n(x_n)]$.
Thus we have that both
$\Gamma_n \Proves \lexists[x_n][!A_n(x_n)]$ and
$\Gamma_n \Proves \lforall[x_n][\lnot !A_n(x_n)]$,
so $\Gamma_n$ itself is inconsistent.
(Note that
$\lforall[x_n][\lnot !A_n(x_n)] \Proves
  \lnot\lexists[x_n][!A_n(x_n)]$.)
Contradiction: $\Gamma_n$ was supposed to be consistent.  Hence
$\Gamma_n \cup \{ !D_n\}$ is consistent.
\end{proof}

\begin{explain}
We'll now show that \emph{complete}, consistent sets which are
saturated have the property that it contains a
  universally quantified !!{sentence} iff it contains all its
  instances and it contains an
  existentially quantified !!{sentence} iff it contains at least one
  instance. We'll use this to show that the !!{structure} we'll
generate from a complete, consistent, saturated set makes all its
quantified sentences true.
\end{explain}

\begin{prop}\ollabel{prop:saturated-instances}
Suppose $\Gamma$ is complete, consistent, and saturated.
\begin{tagenumerate}{prvEx,prvAll}
$\lexists[x][!A(x)] \in \Gamma$ iff $!A(t) \in \Gamma$
  for at least one closed term~$t$.
$\lforall[x][!A(x)] \in \Gamma$ iff $!A(t) \in \Gamma$
  for all closed terms~$t$.
\end{tagenumerate}
\end{prop}

\begin{proof}
  \begin{tagenumerate}{prvEx,prvAll}
  
    First suppose that $\lexists[x][!A(x)]
      \in \Gamma$.  Because $\Gamma$ is saturated,
      $(\lexists[x][!A(x)] \lif !A(c)) \in \Gamma$ for some
      !!{constant}~$c$. By
      \tagrefs{prfAX/{fol:axd:ppr:prop:provability-lif},prfSC/{fol:seq:ppr:prop:provability-lif},prfND/{fol:ntd:ppr:prop:provability-lif},prfTab/{fol:tab:ppr:prop:provability-lif}}, item~(1),
      and \olref[ccs]{prop:ccs}\olref[ccs]{prop:ccs-prov-in}, $!A(c)
      \in \Gamma$.

    For the other direction, saturation is not necessary: Suppose
    $!A(t) \in \Gamma$. Then $\Gamma \Proves \lexists[x][!A(x)]$ by
      \tagrefs{prfAX/{fol:axd:qpr:prop:provability-quantifiers},prfSC/{fol:seq:qpr:prop:provability-quantifiers},prfND/{fol:ntd:qpr:prop:provability-quantifiers},prfTab/{fol:tab:qpr:prop:provability-quantifiers}}, item~(1). By
    \olref[ccs]{prop:ccs}\olref[ccs]{prop:ccs-prov-in},
    $\lexists[x][!A(x)] \in \Gamma$.
  
    Exercise.
  \end{tagenumerate}
\end{proof}






\olfileid{fol}{com}{lin}

\olsection{Lindenbaum's Lemma}

\begin{explain}
We now prove a lemma that shows that any consistent set of
!!{sentence}s is contained in some set of sentences which is not just
consistent, but also !!{complete}. The proof works by adding one
!!{sentence} at a time, guaranteeing at each step that the set remains
consistent. We do this so that for every $!A$, either $!A$ or $\lnot
!A$ gets added at some stage. The union of all stages in that
construction then contains either $!A$ or its negation~$\lnot !A$ and
is thus complete. It is also consistent, since we make sure at each
stage not to introduce an inconsistency.
\end{explain}

\begin{lem}[Lindenbaum's Lemma]
\ollabel{lem:lindenbaum} Every consistent
set~$\Gamma$ in a language~$\Lang{L}$ can be
extended to !!a{complete} and consistent set~$\Gamma^*$.
\end{lem}

\begin{proof}
Let $\Gamma$ be consistent.  Let $!A_0$, $!A_1$,
\dots{} be an enumeration of all the !!{sentence}s of~$\Lang L$.
Define $\Gamma_0 = \Gamma$, and
\[
\Gamma_{n+1} =
\begin{cases}
\Gamma_n \cup \{ !A_n \} & \textrm{if $\Gamma_n \cup \{!A_n\}$ is
  consistent;} \\
\Gamma_n \cup \{ \lnot !A_n \} & \textrm{otherwise.}
\end{cases}
\]
Let $\Gamma^* = \bigcup_{n \geq 0} \Gamma_n$.

Each $\Gamma_n$ is consistent: $\Gamma_0$ is consistent by definition.
If $\Gamma_{n+1} = \Gamma_n \cup \{!A_n\}$, this is because the latter
is consistent.  If it isn't, $\Gamma_{n+1} = \Gamma_n \cup \{\lnot
!A_n\}$. We have to verify that $\Gamma_n \cup \{\lnot !A_n\}$ is
consistent. Suppose it's not. Then \emph{both} $\Gamma_n \cup
\{!A_n\}$ and $\Gamma_n \cup \{\lnot !A_n\}$ are inconsistent.  This
means that $\Gamma_n$ would be inconsistent by

  \tagrefs{prfAX/{fol:axd:prv:prop:provability-exhaustive},
    prfSC/{fol:seq:prv:prop:provability-exhaustive},
    prfND/{fol:ntd:prv:prop:provability-exhaustive},
    prfTab/{fol:tab:prv:prop:provability-exhaustive}},
contrary to the induction hypothesis.

For every~$n$ and every $i < n$, $\Gamma_i \subseteq \Gamma_n$. This
follows by a simple induction on~$n$. For $n=0$, there are no $i < 0$,
so the claim holds automatically.  For the inductive step, suppose it
is true for~$n$. We show that if $i < n+1$ then $\Gamma_i \subseteq
\Gamma_{n+1}$. We have $\Gamma_{n+1} = \Gamma_n \cup \{!A_n\}$ or $=
\Gamma_n \cup \{\lnot !A_n\}$ by construction. So $\Gamma_n \subseteq
\Gamma_{n+1}$. If $i < n+1$, then $\Gamma_i \subseteq \Gamma_n$ by
inductive hypothesis (if $i < n$) or the trivial fact that $\Gamma_n
\subseteq \Gamma_n$ (if $i = n$). We get that $\Gamma_i \subseteq
\Gamma_{n+1}$ by transitivity of~$\subseteq$.

From this it follows that $\Gamma^*$ is consistent. Here's why: Let
$\Gamma' \subseteq \Gamma^*$ be finite. Each $!B \in \Gamma'$ is also
in~$\Gamma_i$ for some~$i$. Let $n$ be the largest of these. Since
$\Gamma_i \subseteq \Gamma_n$ if $i \le n$, every $!B \in \Gamma'$ is
also $\in \Gamma_n$, i.e., $\Gamma' \subseteq \Gamma_n$, and
$\Gamma_n$~is consistent. So, every finite subset $\Gamma' \subseteq
\Gamma^*$ is consistent. By 
  \tagrefs{prfAX/{fol:axd:ptn:prop:proves-compact},
    prfSC/{fol:seq:ptn:prop:proves-compact},
    prfND/{fol:ntd:ptn:prop:proves-compact},
    prfTab/{fol:tab:ptn:prop:proves-compact}}, $\Gamma^*$ is
  consistent.

Every !!{sentence} of $\Frm[L]$ appears on the list used to
define~$\Gamma^*$. If $!A_n \notin \Gamma^*$, then that is because
$\Gamma_n \cup \{!A_n\}$ was inconsistent.  But then $\lnot !A_n
\in \Gamma^*$, so $\Gamma^*$ is !!{complete}.
\end{proof}





\olfileid{fol}{com}{mod}

\olsection{Construction of a Model}

\begin{explain}
Right now we are not concerned about $\eq$, i.e., we only
  want to show that a consistent set~$\Gamma$ of !!{sentence}s not
  containing~$\eq$ is satisfiable. We first extend~$\Gamma$ to a
  consistent, !!{complete}, and saturated set~$\Gamma^*$.  In this
  case, the definition of a model~$\Struct{M(\Gamma^*)}$ is simple: We
  take the set of closed terms of~$\Lang{L'}$ as the domain. We assign
  every !!{constant} to itself, and make sure that more generally, for
  every closed term~$t$, $\Value{t}{M(\Gamma^*)} = t$.  The
  !!{predicate}s are assigned extensions in such a way that an atomic
  !!{sentence} is true in $\Struct{M(\Gamma^*)}$ iff it is
  in~$\Gamma^*$.  This will obviously make all the atomic
  !!{sentence}s in~$\Gamma^*$ true in $\Struct{M(\Gamma^*)}$. The rest
  are true provided the $\Gamma^*$ we start with is consistent,
  complete, and saturated.
\end{explain}


\begin{defn}[Term model]
\ollabel{defn:termmodel}
Let $\Gamma^*$ be !!a{complete} and consistent,
saturated set of !!{sentence}s in a language~$\Lang L$. The \emph{term
  model}~$\Struct M(\Gamma^*)$ of $\Gamma^*$ is the !!{structure}
defined as follows:
\begin{enumerate}
\item The !!{domain}~$\Domain{M(\Gamma^*)}$ is the set of all closed
  terms of~$\Lang L$.
\item The interpretation of !!a{constant} $c$ is $c$ itself:
  $\Assign{c}{M(\Gamma^*)} = c$.
\item The !!{function}~$f$ is assigned the function which, given as
  arguments the closed terms $t_1$, \dots, $t_n$, has as value the
  closed term $f(t_1, \dots, t_n)$:
\[
\Assign{f}{M(\Gamma^*)}(t_1, \dots, t_n) = f(t_1,\dots, t_n)
\]
\item If $R$ is an $n$-place !!{predicate}, then
  \[
  \tuple{t_1, \dots,
  t_n} \in \Assign{R}{M(\Gamma^*)} \text{ iff } \Atom{R}{t_1, \dots,
    t_n} \in \Gamma^*.
  \]
\end{enumerate}
\end{defn}



We will now check that we indeed have $\Value{t}{M(\Gamma^*)} = t$.

\begin{lem} 
\ollabel{lem:val-in-termmodel} Let $\Struct M(\Gamma^*)$ be the term model
of \olref{defn:termmodel}, then $\Value{t}{M(\Gamma^*)} = t$.
\end{lem}
\begin{proof} 
 The proof is by induction on $t$, where the base case, when $t$
 is !!a{constant}, follows directly from the definition of the term
 model. For the induction step assume $t_1, \ldots, t_n$ are closed terms
 such that $\Value{t_i}{M(\Gamma^*)} = t_i$ and that $f$ is an $n$-ary
 !!{function}. Then 
\begin{align*}
\Value{f(t_1,\ldots,t_n)}{M(\Gamma^*)} &= \Assign{f}{M(\Gamma^*)}(\Value{t_1}
 {M(\Gamma^*)}, 
\ldots, \Value{t_n}{M(\Gamma^*)}) \\
&= \Assign{f}{M(\Gamma^*)}(t_1, \dots, t_n) \\
&= f(t_1,\dots, t_n),
\end{align*} 
and so by induction this holds for every closed term~$t$.
\end{proof}
{FOL}{
\begin{explain}
  !!^a{structure}~$\Struct{M}$ may make an existentially
    quantified !!{sentence}~$\lexists[x][!A(x)]$ true without there
    being an instance~$!A(t)$ that it makes true.
  !!^a{structure}~$\Struct{M}$ may make all
    instances~$!A(t)$ of a universally quantified
    !!{sentence}~$\lforall[x][!A(x)]$ true, without
    making~$\lforall[x][!A(x)]$ true. This is because in general
  not every !!{element} of~$\Domain{M}$ is the value of a closed term
  ($\Struct{M}$ may not be covered).  This is the reason the
  satisfaction relation is defined via variable assignments. However,
  for our term model~$\Struct{M(\Gamma^*)}$ this wouldn't be
  necessary---because it is covered. This is the content of the next
  result.
\end{explain}

\begin{prop}
\ollabel{prop:quant-termmodel}
Let $\Struct M(\Gamma^*)$ be the term model of \olref{defn:termmodel}.
\begin{tagenumerate}{prvEx,prvAll}
$\Sat{M(\Gamma^*)}{\lexists[x][!A(x)]}$ iff
  $\Sat{M(\Gamma^*)}{!A(t)}$ for at least one closed term~$t$.
$\Sat{M(\Gamma^*)}{\lforall[x][!A(x)]}$ iff
  $\Sat{M(\Gamma^*)}{!A(t)}$ for all closed terms~$t$.
\end{tagenumerate}
\end{prop}

\begin{proof}
\begin{tagenumerate}{prvEx,prvAll}

  By \olref[syn][ass]{prop:sat-quant},
    $\Sat{M(\Gamma^*)}{\lexists[x][!A(x)]}$ iff for at least one
    variable assignment~$s$, $\Sat{M(\Gamma^*)}{!A(x)}[s]$. As
    $\Domain{M(\Gamma^*)}$ consists of the closed terms of~$\Lang{L}$,
    this is the case iff there is at least one closed term~$t$ such
    that $s(x) = t$ and $\Sat{M(\Gamma^*)}{!A(x)}[s]$.  By
    \olref[fol][syn][ext]{prop:ext-formulas},
    $\Sat{M(\Gamma^*)}{!A(x)}[s]$ iff $\Sat{M(\Gamma^*)}{!A(t)}[s]$,
    where $s(x) = t$.  By \olref[fol][syn][ass]{prop:sentence-sat-true},
    $\Sat{M(\Gamma^*)}{!A(t)}[s]$ iff $\Sat{M(\Gamma^*)}{!A(t)}$,
    since $!A(t)$ is a sentence.

  Exercise.
\end{tagenumerate}
\end{proof}
}{}


\begin{prob}
Complete the proof of \olref[fol][com][mod]{prop:quant-termmodel}.
\end{prob}


\begin{lem}[Truth Lemma]
\ollabel{lem:truth} Suppose $!A$ does not contain~$\eq$. Then
$\Sat{M(\Gamma^*)}{!A}$ iff $!A \in \Gamma^*$.
\end{lem}

\begin{proof}
We prove both directions simultaneously, and by induction on $!A$.
\begin{enumerate}
\indcase{!A}{\lfalse}
  {$\Sat/{M(\Gamma^*)}{\lfalse}$
    by definition of satisfaction. On the other hand, $\lfalse \notin
    \Gamma^*$ since $\Gamma^*$ is consistent.}



\indcase{!A}{R(t_1, \dots, t_n)}
  {$\Sat{M(\Gamma^*)}{\Atom{R}{t_1, \dots, t_n}}$ iff $\tuple{t_1,
      \dots, t_n} \in \Assign{R}{M(\Gamma^*)}$ (by the definition of
    satisfaction) iff $R(t_1, \dots, t_n) \in \Gamma^*$ (by the
    construction of $\Struct
    M(\Gamma^*)$).}


  \indcase{!A}{\lnot !B}
    {$\Sat{M(\Gamma^*)}{\indfrm}$ iff
    $\Sat/{M(\Gamma^*)}{!B}$ (by
    definition of satisfaction). By induction hypothesis,
    $\Sat/{M(\Gamma^*)}{!B}$ iff
    $!B \notin \Gamma^*$. Since $\Gamma^*$ is consistent and
    !!{complete}, $!B
    \notin \Gamma^*$ iff $\lnot !B \in \Gamma^*$.}



  \indcase!{!A}{!B \land !C}{}



  \indcase{!A}{!B \lor !C}
    {$\Sat{M(\Gamma^*)}{\indfrm}$
    iff $\Sat{M(\Gamma^*)}{!B}$ or
    $\Sat{M(\Gamma^*)}{!C}$ (by
    definition of satisfaction) iff $!B \in \Gamma^*$ or $!C \in
    \Gamma^*$ (by induction hypothesis). This is the case iff $(!B
    \lor !C) \in \Gamma^*$ (by
    \olref[ccs]{prop:ccs}\olref[ccs]{prop:ccs-or}).}
  


  \indcase!{!A}{!B \lif !C}{}



  
    \indcase!{!A}{\lforall[x][!B(x)]}{}


  
    \indcase{!A}{\lexists[x][!B(x)]}{$\Sat{M(\Gamma^*)}{\indfrm}$ iff
      $\Sat{M(\Gamma^*)}{!B(t)}$ for at least one term~$t$
      (\olref{prop:quant-termmodel}).  By induction hypothesis, this
      is the case iff $!B(t) \in \Gamma^*$ for at least one term~$t$.
      By \olref[hen]{prop:saturated-instances}, this in turn is the
      case iff $\lexists[x][!B(x)] \in \Gamma^*$.}

\end{enumerate}
\end{proof}




\begin{prob}
Complete the proof of \olref[fol][com][mod]{lem:truth}.
\end{prob}








  

\olfileid{fol}{com}{ide}
\olsection{Identity}

\begin{explain}
The construction of the term model given in the preceding section is
enough to establish completeness for first-order logic for
sets~$\Gamma$ that do not contain~$\eq$.  The term model satisfies
every $!A \in \Gamma^*$ which does not contain~$\eq$ (and hence all
$!A \in \Gamma$).  It does not work, however, if $\eq$ is present.
The reason is that $\Gamma^*$ then may contain
!!a{sentence}~$\eq[t][t']$, but in the term model the value of any
term is that term itself.  Hence, if $t$ and $t'$ are different terms,
their values in the term model---i.e., $t$ and $t'$,
respectively---are different, and so $\eq[t][t']$ is false.  We can
fix this, however, using a construction known as ``factoring.''
\end{explain}

\begin{defn}
  Let $\Gamma^*$ be a consistent and !!{complete} set of sentences
  in~$\Lang L$.  We define the relation $\approx$ on the set of closed
  terms of~$\Lang L$ by
  \[
  t \approx t' \text{\quad iff \quad} \eq[t][t'] \in \Gamma^*
  \]
\end{defn}

\begin{prop}
\ollabel{prop:approx-equiv}
The relation $\approx$ has the following properties:
\begin{enumerate}
\item $\approx$ is reflexive.
\item $\approx$ is symmetric.
\item  $\approx$ is transitive.
\item If $t \approx t'$, $f$ is !!a{function}, and $t_1$, \dots,
  $t_{i-1}$, $t_{i+1}$, \dots, $t_n$ are closed terms, then
\[
\Atom{f}{t_1,\dots, t_{i-1}, t, t_{i+1}, \dots, t_n} \approx
\Atom{f}{t_1,\dots, t_{i-1}, t', t_{i+1}, \dots, t_n}.
\]
\item If $t \approx t'$, $R$ is !!a{predicate}, and $t_1$, \dots,
  $t_{i-1}$, $t_{i+1}$, \dots, $t_n$ are closed terms, then
\begin{multline*}
\Atom{R}{t_1,\dots, t_{i-1}, t, t_{i+1}, \dots, t_n} \in \Gamma^* \text{ iff } \\
\Atom{R}{t_1,\dots, t_{i-1}, t', t_{i+1}, \dots, t_n} \in \Gamma^*.
\end{multline*}
\end{enumerate}
\end{prop}

\begin{proof}
Since $\Gamma^*$ is consistent and !!{complete}, $\eq[t][t'] \in
\Gamma^*$ iff $\Gamma^* \Proves \eq[t][t']$.  Thus it is enough to
show the following:
\begin{enumerate}
\item $\Gamma^* \Proves \eq[t][t]$ for all closed terms~$t$.
\item If $\Gamma^* \Proves \eq[t][t']$ then $\Gamma^* \Proves \eq[t'][t]$.
\item If $\Gamma^* \Proves \eq[t][t']$ and $\Gamma^* \Proves
  \eq[t'][t'']$, then $\Gamma^* \Proves \eq[t][t'']$.
\item If $\Gamma^* \Proves \eq[t][t']$, then
\[
\Gamma^* \Proves
\eq[\Atom{f}{t_1,\dots,t_{i-1},t,t_{i+1},,\dots,t_n}][\Atom{f}{t_1,\dots,t_{i-1},t',t_{i+1},\dots,t_n}]
\]
for every $n$-place !!{function}~$f$ and closed terms $t_1$, \dots,
$t_{i-1}$, $t_{i+1}$, \dots,~$t_n$.
\item If $\Gamma^* \Proves \eq[t][t']$ and
$\Gamma^* \Proves
\Atom{R}{t_1,\dots,t_{i-1},t,t_{i+1},\dots,t_n}$, then
$\Gamma^* \Proves \Atom{R}{t_1,\dots,t_{i-1},t',t_{i+1},\dots,t_n}$
for every $n$-place !!{predicate}~$R$ and closed terms $t_1$, \dots,
$t_{i-1}$, $t_{i+1}$, \dots,~$t_n$.
\end{enumerate}
\end{proof}

\begin{prob}
Complete the proof of~\olref[fol][com][ide]{prop:approx-equiv}.
\end{prob}

\begin{defn}
Suppose $\Gamma^*$ is a consistent and !!{complete} set in a
language~$\Lang L$, $t$ is a closed term, and $\approx$ as in the
previous definition. Then:
\[
\equivrep{t}{\approx} = \Setabs{t'}{t'\in \Trm[L], t \approx t'}
\]
and $\equivclass{\Trm[L]}{\approx} = \Setabs{\equivrep{t}{\approx}}{t \in \Trm[L]}$.
\end{defn}

\begin{defn}
\ollabel{defn:term-model-factor}
Let $\Struct M = \Struct M(\Gamma^*)$ be the term model
for~$\Gamma^*$ from \olref[mod]{defn:termmodel}.  Then $\Struct{\equivclass{M}{\approx}}$ is the following
!!{structure}:
\begin{enumerate}
\item $\Domain{\equivclass{M}{\approx}} = \equivclass{\Trm[L]}{\approx}$.
\item $\Assign{c}{\equivclass{M}{\approx}} = \equivrep{c}{\approx}$
\item $\Assign{f}{\equivclass{M}{\approx}}(\equivrep{t_1}{\approx}, \dots,
  \equivrep{t_n}{\approx}) = \equivrep{\Atom{f}{t_1,\dots, t_n}}{\approx}$
\item $\tuple{\equivrep{t_1}{\approx}, \dots, \equivrep{t_n}{\approx}} \in
  \Assign{R}{\equivclass{M}{\approx}}$ iff
  $\Sat{M}{\Atom{R}{t_1,\dots, t_n}}$, i.e., iff $\Atom{R}{t_1,\dots,
  t_n} \in \Gamma^*$.
\end{enumerate}
\end{defn}

\begin{explain}
Note that we have defined $\Assign{f}{\equivclass{M}{\approx}}$ and
$\Assign{R}{\equivclass{M}{\approx}}$ for elements of
$\equivclass{\Trm[L]}{\approx}$ by referring to them as
$\equivrep{t}{\approx}$, i.e., via \emph{representatives}~$t \in
\equivrep{t}{\approx}$.  We have to make sure that these definitions
do not depend on the choice of these representatives, i.e., that for
some other choices~$t'$ which determine the same equivalence classes
($\equivrep{t}{\approx} = \equivrep{t'}{\approx}$), the definitions
yield the same result. For instance, if $R$ is a one-place
!!{predicate}, the last clause of the definition says that
$\equivrep{t}{\approx} \in \Assign{R}{\equivclass{M}{\approx}}$ iff
$\Sat{M}{\Atom{R}{t}}$. If for some other term~$t'$ with $t \approx
t'$, $\Sat/{M}{\Atom{R}{t}}$, then the definition would require
$\equivrep{t'}{\approx} \notin \Assign{R}{\equivclass{M}{\approx}}$.
If $t \approx t'$, then $\equivrep{t}{\approx} =
\equivrep{t'}{\approx}$, but we can't have both $\equivrep{t}{\approx}
\in \Assign{R}{\equivclass{M}{\approx}}$ and $\equivrep{t}{\approx}
\notin \Assign{R}{\equivclass{M}{\approx}}$.  However,
\olref{prop:approx-equiv} guarantees that this cannot happen.
\end{explain}

\begin{prop}
$\Struct{\equivclass{M}{\approx}}$ is well defined, i.e., if $t_1$,
  \dots, $t_n$, $t_1'$, \dots, $t_n'$ are closed terms,
  and $t_i \approx t_i'$ then
\begin{enumerate}
\item $\equivrep{\Atom{f}{t_1,\dots, t_n}}{\approx} =
    \equivrep{\Atom{f}{t_1',\dots, t_n'}}{\approx}$, i.e.,
  \[
  \Atom{f}{t_1,\dots, t_n} \approx \Atom{f}{t_1',\dots, t_n'}
  \]
  and
\item $\Sat{M}{\Atom{R}{t_1,\dots, t_n}}$ iff
  $\Sat{M}{\Atom{R}{t_1',\dots, t_n'}}$, i.e.,
  \[
    \Atom{R}{t_1,\dots, t_n} \in \Gamma^* \text{ iff }
    \Atom{R}{t_1',\dots, t_n'} \in \Gamma^*.
  \]
\end{enumerate}
\end{prop}

\begin{proof}
Follows from \olref{prop:approx-equiv} by induction on~$n$.
\end{proof}

As in the case of the term model, before proving the truth lemma we need the
following lemma.

\begin{lem} 
\ollabel{lem:val-in-termmodel-factored} Let $\Struct M = \Struct M
 (\Gamma^*)$, then $\Value{t}{\equivclass{M}{\approx}} = \equivrep{t}
 {\approx}$.
\end{lem}
\begin{proof} 
The proof is similar to that of \olref[mod]{lem:val-in-termmodel}.
\end{proof}

\begin{prob}
Complete the proof of~\olref[fol][com][ide]{lem:val-in-termmodel-factored}.
\end{prob}

\begin{lem}
\ollabel{lem:truth}
$\Sat{\equivclass{M}{\approx}}{!A}$ iff $!A \in \Gamma^*$ for all
  sentences~$!A$.
\end{lem}

\begin{proof}
By induction on~$!A$, just as in the proof of \olref[mod]{lem:truth}.
The only case that needs additional attention is when $!A \ident
\eq[t][t']$.
\begin{align*}
\Sat{\equivclass{M}{\approx}}{\eq[t][t']} & \text{ iff } \equivrep{t}{\approx} = \equivrep{t'}{\approx}
\text{ (by definition of $\Struct{\equivclass{M}{\approx}}$)}\\
& \text{ iff } t \approx t' \text{ (by definition of $\equivrep{t}{\approx}$)}\\
& \text{ iff } \eq[t][t'] \in \Gamma^* \text{ (by definition of $\approx$).}
\end{align*}
\end{proof}

\begin{digress}
Note that while $\Struct{M(\Gamma^*)}$ is always !!{enumerable} and
infinite, $\Struct{\equivclass{M}{\approx}}$ may be finite, since it
may turn out that there are only finitely many classes
$\equivrep{t}{\approx}$.  This is to be expected, since $\Gamma$ may
contain !!{sentence}s which require any !!{structure} in which they
are true to be finite.  For instance, $\lforall[x][\lforall[y][x =
y]]$ is a consistent !!{sentence}, but is satisfied only in
!!{structure}s with !!a{domain} that contains exactly one
!!{element}.
\end{digress}






\olfileid{fol}{com}{cth}

\olsection{The Completeness Theorem}

\begin{explain}
Let's combine our results: we arrive at the completeness theorem.
\end{explain}

\begin{thm}[Completeness Theorem]
\ollabel{thm:completeness}
Let $\Gamma$ be a set of !!{sentence}s.  If $\Gamma$ is consistent, it
is satisfiable.
\end{thm}

\begin{proof}
Suppose $\Gamma$ is consistent. By
  \olref[hen]{lem:henkin}, there is a saturated consistent set
  $\Gamma' \supseteq \Gamma$. By \olref[lin]{lem:lindenbaum}, there
is a $\Gamma^* \supseteq \Gamma'$ which is
consistent and !!{complete}.

  Since $\Gamma' \subseteq \Gamma^*$, for each
  !!{formula}~$!A(x)$, $\Gamma^*$ contains !!a{sentence} of the form
  $\lexists[x][!A(x)] \lif !A(c)$
  and so $\Gamma^*$ is saturated.  If $\Gamma$
  does not contain~$\eq$, then by
\olref[mod]{lem:truth},
$\Sat{M(\Gamma^*)}{!A}$ iff $!A \in
\Gamma^*$.  From this it follows in particular that for all $!A \in
\Gamma$, $\Sat{M(\Gamma^*)}{!A}$, so
$\Gamma$ is satisfiable.
  If $\Gamma$ does contain~$\eq$,
  then by \olref[ide]{lem:truth}, for all !!{sentence}s~$!A$,
  $\Sat{\equivclass{M}{\approx}}{!A}$ iff $!A \in \Gamma^*$.  In
  particular, $\Sat{\equivclass{M}{\approx}}{!A}$ for all $!A \in
  \Gamma$, so $\Gamma$ is satisfiable.
\end{proof}

\begin{cor}[Completeness Theorem, Second Version]
\ollabel{cor:completeness}
For all $\Gamma$ and !!{sentence}s~$!A$: if $\Gamma \Entails !A$ then
$\Gamma \Proves !A$.
\end{cor}

\begin{proof}
Note that the $\Gamma$'s in \olref{cor:completeness} and
\olref{thm:completeness} are universally quantified.  To make sure we
do not confuse ourselves, let us restate \olref{thm:completeness}
using a different variable: for any set of !!{sentence}s~$\Delta$, if
$\Delta$ is consistent, it is satisfiable.  By contraposition, if
$\Delta$ is not satisfiable, then $\Delta$ is inconsistent.  We will
use this to prove the corollary.

Suppose that $\Gamma \Entails !A$.  Then $\Gamma \cup \{\lnot !A\}$ is
unsatisfiable by \olref[syn][sem]{prop:entails-unsat}.  Taking $\Gamma
\cup \{\lnot !A\}$ as our $\Delta$, the previous version of
\olref{thm:completeness} gives us that $\Gamma \cup \{\lnot !A\}$ is
inconsistent.  By

  \tagrefs{prfAX/{fol:axd:prv:prop:prov-incons},
    prfSC/{fol:seq:prv:prop:prov-incons},
    prfND/{fol:ntd:prv:prop:prov-incons},
    prfTab/{fol:tab:prv:prop:prov-incons}},
$\Gamma \Proves !A$.
\end{proof}


\begin{prob}
Use \olref[fol][com][cth]{cor:completeness} to prove
\olref[fol][com][cth]{thm:completeness}, thus showing that the two
formulations of the completeness theorem are equivalent.
\end{prob}





\begin{prob}
In order for !!a{derivation} system to be complete, its rules must be
strong enough to prove every unsatisfiable set inconsistent.  Which of
the rules of !!{derivation} were necessary to prove completeness?  Are any
of these rules not used anywhere in the proof?  In order to answer
these questions, make a list or diagram that shows which of the rules
of !!{derivation} were used in which results that lead up to the proof of
\olref[fol][com][cth]{thm:completeness}.  Be sure to note any tacit
uses of rules in these proofs.
\end{prob}








\olfileid{fol}{com}{com}

\olsection{The Compactness Theorem}

One important consequence of the completeness theorem is the
compactness theorem.  The compactness theorem states that if each
\emph{finite} subset of a set of !!{sentence}s is satisfiable, the
entire set is satisfiable---even if the set itself is infinite. This
is far from obvious. There is nothing that seems to rule out, at first
glance at least, the possibility of there being infinite sets of
!!{sentence}s which are contradictory, but the contradiction only
arises, so to speak, from the infinite number.  The compactness
theorem says that such a scenario can be ruled out: there are no
unsatisfiable infinite sets of !!{sentence}s each finite subset of
which is satisfiable. Like the completeness theorem, it has a version
related to entailment: if an infinite set of !!{sentence}s entails
something, already a finite subset does.

\begin{defn}
  A set $\Gamma$ of !!{formula}s is \emph{finitely satisfiable} iff every finite $\Gamma_0 \subseteq \Gamma$ is satisfiable.
\end{defn}

\begin{thm}[Compactness Theorem]
\ollabel{thm:compactness}
The following hold for any sentences $\Gamma$ and $!A$:
\begin{enumerate}
  \item $\Gamma \Entails !A$ iff there is a finite $\Gamma_0
    \subseteq \Gamma$ such that $\Gamma_0 \Entails !A$.
  \item $\Gamma$ is satisfiable iff it is finitely
    satisfiable.
\end{enumerate}
\end{thm}

\begin{proof}
We prove (2).  If $\Gamma$ is satisfiable, then there is
!!a{structure}~$\Struct{M}$
such that $\Sat{M}{!A}$ for all $!A \in
\Gamma$.  Of course, this $\Struct{M}$ also
satisfies every finite subset of~$\Gamma$, so $\Gamma$ is finitely
satisfiable.

Now suppose that $\Gamma$ is finitely satisfiable.  Then every finite
subset~$\Gamma_0 \subseteq \Gamma$ is satisfiable.  By soundness
(
  \tagrefs{prfAX/{fol:axd:sou:cor:consistency-soundness},
    prfSC/{fol:seq:sou:cor:consistency-soundness},
    prfND/{fol:ntd:sou:cor:consistency-soundness},
    prfTab/{fol:tab:sou:cor:consistency-soundness}}),
every finite subset is consistent.  Then $\Gamma$ itself must be
consistent by

  \tagrefs{prfAX/{fol:axd:ptn:prop:proves-compact},
    prfSC/{fol:seq:ptn:prop:proves-compact},
    prfND/{fol:ntd:ptn:prop:proves-compact},
    prfTab/{fol:tab:ptn:prop:proves-compact}}.
By completeness (\olref[cth]{thm:completeness}), since $\Gamma$~is
consistent, it is satisfiable.
\end{proof}


\begin{prob}
Prove (1) of \olref[fol][com][com]{thm:compactness}.
\end{prob}





\begin{ex}
In every model~$\Struct{M}$ of a theory~$\Gamma$, each term~$t$ of
course picks out !!a{element} of~$\Domain{M}$. Can we guarantee that
it is also true that every !!{element} of~$\Domain{M}$ is picked out
by some term or other? In other words, are there theories~$\Gamma$ all
models of which are covered?  The compactness theorem shows that this
is not the case if $\Gamma$ has infinite models. Here's how to see
this: Let $\Struct{M}$ be an infinite model of~$\Gamma$, and let $c$
be !!a{constant} not in the language of~$\Gamma$. Let $\Delta$ be the
set of all sentences $\eq/[c][t]$ for $t$ a term in the
language~$\Lang{L}$ of~$\Gamma$, i.e.,
  \[
  \Delta = \Setabs{\eq/[c][t]}{t \in \Trm[L]}.
  \]
A finite subset of $\Gamma \cup \Delta$ can be written as $\Gamma'
\cup \Delta'$, with $\Gamma' \subseteq \Gamma$ and $\Delta' \subseteq
\Delta$. Since $\Delta'$ is finite, it can contain only finitely many
terms. Let $a \in \Domain{M}$ be !!a{element} of $\Domain{M}$ not
picked out by any of them, and let $\Struct{M'}$ be the !!{structure}
that is just like $\Struct{M}$, but also $\Assign{c}{M'} = a$. Since
$a \neq \Value{t}{M}$ for all~$t$ occurring in~$\Delta'$,
$\Sat{M'}{\Delta'}$. Since $\Sat{M}{\Gamma}$, $\Gamma' \subseteq
\Gamma$, and $c$ does not occur in~$\Gamma$, also
$\Sat{M'}{\Gamma'}$. Together, $\Sat{M'}{\Gamma' \cup \Delta'}$ for
every finite subset $\Gamma' \cup \Delta'$ of $\Gamma \cup \Delta$. So
every finite subset of $\Gamma \cup \Delta$ is satisfiable. By
compactness, $\Gamma \cup \Delta$ itself is satisfiable. So there are
models~$\Sat{M}{\Gamma \cup \Delta}$. Every such $\Struct{M}$ is a
model of~$\Gamma$, but is not covered, since $\Value{c}{M} \neq
\Value{t}{M}$ for all terms~$t$ of~$\Lang{L}$.
\end{ex}

\begin{ex}
Consider !!a{language} $\Lang{L}$ containing the !!{predicate}~$<$,
!!{constant}s $\Obj{0}$, $\Obj{1}$, and !!{function}s $+$, $\times$, and
$-$. Let $\Gamma$ be the set of all !!{sentence}s in this
!!{language} true in the !!{structure}~$\Struct{Q}$ with domain~$\Rat$ and the obvious
interpretations.  $\Gamma$~is the set of all !!{sentence}s
of~$\Lang{L}$ true about the rational numbers. Of course, in~$\Rat$
(and even in~$\Real$), there are no numbers~$r$ which are greater than~$0$
but less than $1/k$ for all $k \in \PosInt$.  Such a number, if it
existed, would be an \emph{infinitesimal:} non-zero, but infinitely
small.  The compactness theorem can be used to show that there are 
models of~$\Gamma$ in which infinitesimals exist. We do not have 
!!a{function} for division in our language (division by zero is 
undefined, and !!{function}s have to be interpreted by total functions).
However, we can still express that $r < 1/k$, since this is the case iff
$r \cdot k < 1$. Now let $c$ be a new !!{constant} and let $\Delta$ be 
\[
\{\Obj{0} < c\}
\cup \Setabs{ c \times \num k < \Obj{1} }{k \in \PosInt}
\]
(where $\num{k} = (\Obj{1} + (\Obj{1} + \dots + (\Obj{1} +
\Obj{1})\dots))$ with $k$~$\Obj{1}$'s). For any finite
subset~$\Delta_0$ of~$\Delta$ there is a~$K$ such that for all the
!!{sentence}s $c \times \num{k} < \Obj{1}$ in~$\Delta_0$ have $k < K$.
If we expand $\Struct{Q}$ to~$\Struct{Q'}$ with $\Assign{c}{Q'} = 1/K$
we have that $\Sat{Q'}{\Gamma_0 \cup \Delta_0}$ for any finite
$\Gamma_0 \subseteq \Gamma$, and so $\Gamma \cup \Delta$ is finitely
satisfiable (Exercise: prove this in detail). By compactness, $\Gamma
\cup \Delta$ is satisfiable. Any model~$\Struct{S}$ of $\Gamma \cup
\Delta$ contains an infinitesimal, namely~$\Assign{c}{S}$.
\end{ex}



\begin{prob}
In the standard model of arithmetic~$\Struct{N}$, there is no
!!{element}~$k \in \Domain{N}$ which satisfies every formula~$\num{n}
< x$ (where $\num{n}$ is $\Obj{0}^{\prime\dots\prime}$ with $n$
$\prime$'s).  Use the compactness theorem to show that the set of
!!{sentence}s in the language of arithmetic which are true in the standard
model of arithmetic $\Struct{N}$ are also true in
!!a{structure}~$\Struct{N'}$ that contains !!a{element} which
\emph{does} satisfy every formula $\num{n} < x$.
\end{prob}



\begin{ex}
We know that first-order logic with !!{identity} can express that the
size of the !!{domain} must have some minimal size: The
sentence~$!A_{\ge n}$ (which says ``there are at least $n$ distinct
objects'') is true only in structures where $\Domain{M}$ has at
least~$n$ objects. So if we take
  \[
  \Delta = \Setabs{!A_{\ge n}}{n \ge 1}
  \]
then any model of $\Delta$ must be infinite. Thus, we can guarantee
that a theory only has infinite models by adding~$\Delta$ to it: the
models of $\Gamma \cup \Delta$ are all and only the infinite models
of~$\Gamma$.

So first-order logic can express infinitude. The compactness theorem
shows that it cannot express finitude, however. For suppose some set
of sentences $\Lambda$ were satisfied in all and only finite
!!{structure}s. Then $\Delta \cup \Lambda$ is finitely
satisfiable. Why? Suppose $\Delta' \cup \Lambda' \subseteq \Delta \cup
\Lambda$ is finite with $\Delta' \subseteq \Delta$ and $\Lambda'
\subseteq \Lambda$. Let $n$ be the largest number such that $!A_{\ge n}
\in \Delta'$. $\Lambda$, being satisfied in all finite structures, has
a model~$\Struct{M}$ with finitely many but $\ge n$ !!{element}s. But
then $\Sat{M}{\Delta' \cup \Lambda'}$.  By compactness, $\Delta \cup
\Lambda$ has an infinite model, contradicting the assumption that
$\Lambda$ is satisfied only in finite !!{structure}s.
\end{ex}






\olfileid{fol}{com}{cpd}

\olsection{A Direct Proof of the Compactness Theorem}

We can prove the Compactness Theorem directly, without appealing to
the Completeness Theorem, using the same ideas as in the proof of the
completeness theorem.  In the proof of the Completeness Theorem we
started with a consistent set~$\Gamma$ of !!{sentence}s, expanded it
to a consistent, saturated, and !!{complete}
set~$\Gamma^*$ of !!{sentence}s, and then showed that in the
term
  model~$\Struct{M(\Gamma^*)}$
constructed from $\Gamma^*$, all !!{sentence}s of~$\Gamma$ are true,
so $\Gamma$ is satisfiable.

We can use the same method to show that a finitely satisfiable set of
sentences is satisfiable. We just have to prove the corresponding
versions of the results leading to the truth lemma where we replace
``consistent'' with ``finitely satisfiable.''

\begin{prop}
\ollabel{prop:fsat-ccs}
Suppose $\Gamma$ is !!{complete} and finitely satisfiable. Then:
\begin{enumerate}
$(!A \land !B) \in \Gamma$
  iff both $!A \in \Gamma$ and $!B \in \Gamma$.

$(!A \lor !B) \in \Gamma$ iff
  either $!A \in \Gamma$ or $!B \in \Gamma$.

$(!A \lif !B) \in \Gamma$ iff
  either $!A \notin \Gamma$ or $!B \in \Gamma$.
\end{enumerate}
\end{prop}


\begin{prob}
Prove \olref[fol][com][cpd]{prop:fsat-ccs}. Avoid the use of $\Proves$.
\end{prob}





\begin{lem}
\ollabel{lem:fsat-henkin} Every finitely satisfiable set~$\Gamma$ can
be extended to a saturated finitely satisfiable set~$\Gamma'$.
\end{lem}



\begin{prob}
Prove \olref[fol][com][cpd]{lem:fsat-henkin}. (Hint: The crucial step
is to show that if $\Gamma_n$ is finitely satisfiable, so is $\Gamma_n
\cup \{!D_n\}$, without any appeal to !!{derivation}s or consistency.)
\end{prob}



\begin{prop}\ollabel{prop:fsat-instances}
Suppose $\Gamma$ is complete, finitely satisfiable, and saturated.
\begin{tagenumerate}{prvEx,prvAll}
$\lexists[x][!A(x)] \in \Gamma$ iff $!A(t) \in \Gamma$
  for at least one closed term~$t$.
$\lforall[x][!A(x)] \in \Gamma$ iff $!A(t) \in \Gamma$
  for all closed terms~$t$.
\end{tagenumerate}
\end{prop}



\begin{prob}
Prove \olref[fol][com][cpd]{prop:fsat-instances}.
\end{prob}


\begin{lem}
\ollabel{lem:fsat-lindenbaum} Every finitely satisfiable set~$\Gamma$
can be extended to !!a{complete} and finitely satisfiable
set~$\Gamma^*$.
\end{lem}


\begin{prob}
Prove \olref[fol][com][cpd]{lem:fsat-lindenbaum}.  (Hint: the crucial
step is to show that if $\Gamma_n$ is finitely satisfiable, then
either $\Gamma_n \cup \{!A_n\}$ or $\Gamma_n \cup \{\lnot !A_n\}$ is
finitely satisfiable.)
\end{prob}




\begin{thm}[Compactness]
\ollabel{thm:compactness-direct} $\Gamma$ is satisfiable if and only
if it is finitely satisfiable.
\end{thm}

\begin{proof}
If $\Gamma$ is satisfiable, then there is
!!a{structure}~$\Struct{M}$
such that $\Sat{M}{!A}$ for all $!A \in \Gamma$.
Of course, this $\Struct{M}$ also
satisfies every finite subset of~$\Gamma$, so $\Gamma$ is finitely
satisfiable.

Now suppose that $\Gamma$ is finitely satisfiable. By
  \olref{lem:fsat-henkin}, there is a finitely satisfiable, saturated
  set $\Gamma' \supseteq \Gamma$. By \olref{lem:fsat-lindenbaum},
$\Gamma'$ can be extended to !!a{complete} and
finitely satisfiable set~$\Gamma^*$, and $\Gamma^*$ is
  still saturated. Construct the term
  model~$\Struct{M(\Gamma^*)}$
as in \olref[mod]{defn:termmodel}. Note that
\olref[mod]{prop:quant-termmodel} did not rely on the fact that
$\Gamma^*$ is consistent (or !!{complete} or saturated, for that
matter), but just on the fact that $\Struct{M(\Gamma^*)}$ is covered.
The proof of the Truth Lemma (\olref[mod]{lem:truth}) goes through if
we replace references to \olref[ccs]{prop:ccs} and
\olref[hen]{prop:saturated-instances} by references to
\olref{prop:fsat-ccs} and \olref{prop:fsat-instances}
\end{proof}


\begin{prob}
Write out the complete proof of the Truth Lemma
(\olref[fol][com][mod]{lem:truth}) in the version required for the
proof of \olref[fol][com][cpd]{thm:compactness-direct}.
\end{prob}







  

\olfileid{fol}{com}{dls}
\olsection{The L\"owenheim--Skolem Theorem}

The L\"owenheim--Skolem Theorem says that if a theory has an infinite
model, then it also has a model that is at most !!{denumerable}. An
immediate consequence of this fact is that first-order logic cannot
express that the size of !!a{structure} is !!{nonenumerable}: any
!!{sentence} or set of !!{sentence}s satisfied in all
!!{nonenumerable} !!{structure}s is also satisfied in some
!!{enumerable} structure.

\begin{thm} 
\ollabel{thm:downward-ls} If $\Gamma$ is consistent then it has
!!a{enumerable} model, i.e., it is satisfiable in !!a{structure}
whose domain is either finite or !!{denumerable}.
\end{thm}

\begin{proof}
If $\Gamma$ is consistent, the !!{structure}~$\Struct M$ delivered by
the proof of the completeness theorem has a domain $\Domain{M}$ that
is no larger than the set of the terms of the language~$\Lang L$. So
$\Struct M$ is at most !!{denumerable}.
\end{proof}

\begin{thm}
\ollabel{noidentity-ls} If $\Gamma$ is a consistent set of !!{sentence}s
in the language of first-order logic without identity, then it has
!!a{denumerable} model, i.e., it is satisfiable in !!a{structure}
whose domain is infinite and !!{enumerable}.
\end{thm}

\begin{proof}
If $\Gamma$ is consistent and contains no sentences in which identity
appears, then the !!{structure}~$\Struct M$ delivered by the proof of
the completeness theorem has a domain $\Domain{M}$ identical to the set
of terms of the language~$\Lang L'$. So $\Struct{M}$ is
!!{denumerable}, since $\Trm[L']$ is.
\end{proof}

\begin{ex}[Skolem's Paradox]
Zermelo--Fraenkel set theory~$\Log{ZFC}$ is a very powerful framework
in which practically all mathematical statements can be expressed,
including facts about the sizes of sets. So for instance, $\Log{ZFC}$
can prove that the set~$\Real$ of real numbers is !!{nonenumerable},
it can prove Cantor's Theorem that the power set of any set is larger
than the set itself, etc.  If $\Log{ZFC}$ is consistent, its models
are all infinite, and moreover, they all contain !!{element}s about
which the theory says that they are !!{nonenumerable}, such as the
element that makes true the theorem of~$\Log{ZFC}$ that the power set
of the natural numbers exists. By the L\"owenheim--Skolem Theorem,
$\Log{ZFC}$ also has !!{enumerable} models---models that contain
``!!{nonenumerable}'' sets but which themselves are !!{enumerable}.
\end{ex}




\OLEndChapterHook





